(1) 在两个δ势之间粒子自由传播，\(k=\sqrt{2mE}/\hbar\)，通解为 \(\psi(x)=Ae^{ikx}+Be^{-ikx}\)。每个 \(x=ja\) 处 \(\psi\) 连续，导数跃变由薛定谔方程小区间积分给出 \(\psi'(ja^+)-\psi'(ja^-)=2mg\,\psi(ja)/\hbar^2\)。这里 \(g\) 的单位是能量乘长度，便于与自由传播矩阵中的 \(k\) 组合成无量纲参数。
(2) 采用状态向量 \((\psi,\psi')^T\)，自由传播距离 \(a\) 的矩阵为 \(P=\begin{pmatrix}\cos ka&\sin ka/k\\-k\sin ka&\cos ka\end{pmatrix}\)，穿越δ势的跳跃矩阵为 \(D=\begin{pmatrix}1&0\\2mg/\hbar^2&1\end{pmatrix}\)。一个晶胞转移矩阵可取 \(M=DP\)，其行列式为一。由 Bloch 条件本征值为 \(e^{\pm iqa}\)，故 \(2\cos(qa)=\operatorname{tr}M=2\cos ka+2mg\sin ka/(\hbar^2k)\)，即 \(\cos(qa)=\cos(ka)+[mg/(\hbar^2k)]\sin(ka)\)。
(3) 实数 Bloch 波数存在的条件是 \(|\cos(ka)+mg\sin(ka)/(\hbar^2k)|\le1\)，否则 \(q\) 具有非零虚部，为禁带。在 \(k\to0\) 时右端趋于 \(1+mga/\hbar^2>1\)，故排斥性周期δ势的低能区存在禁止传播的区间；当 \(g\to0\)，色散回到 \(\cos(qa)=\cos(ka)\)，即自由粒子抛物线按晶格倒格矢折叠。
(4) 在区边界 \(q=\pm\pi/a\)，自由粒子的两个反向传播平面波在折叠表象发生简并，周期势的倒格矢傅里叶分量将它们耦合成驻波、分裂原能量，从而打开带隙。孤立δ势仅造成连续散射振幅并不形成无限周期 Bloch 边界条件；能隙来自大量周期重复散射的相干干涉，而不是单个δ势存在非零透射损耗。
